\documentclass[11pt]{article}
\usepackage{mathblatt}
% vorspann.tex – Kompetenzblatt, zusammenbau v0.7 (2026-09-28).
% Ergänzt mathblatt.sty für Kompetenzblätter, ohne sie zu ändern
% (layout-befunde.md 1–34 vom 28.09.). Wird nach \usepackage{mathblatt}
% eingelesen.
\makeatletter
\RequirePackage{newunicodechar}
% Zeichen, die Latin Modern nicht hat (Befund 20): Text aus einer
% Ersatzschrift, Mathe als Befehl; ohne Ersatzschrift auch Text als Mathe.
\newif\ifkbersatz
\IfFontExistsTF{FreeSerif}{\newfontfamily\kbersatzschrift{FreeSerif}\kbersatztrue}{%
 \IfFontExistsTF{DejaVu Serif}{\newfontfamily\kbersatzschrift{DejaVu Serif}\kbersatztrue}{%
  \IfFontExistsTF{Cambria Math}{\newfontfamily\kbersatzschrift{Cambria Math}\kbersatztrue}{%
   \IfFontExistsTF{Segoe UI Symbol}{\newfontfamily\kbersatzschrift{Segoe UI Symbol}\kbersatztrue}{}}}}
\newcommand{\kbzeichen}[2]{\ifmmode #2\else\ifkbersatz{\kbersatzschrift #1}\else\ensuremath{#2}\fi\fi}
\newcommand{\kbzeichenhoch}[1]{\ifkbersatz\ifmmode\text{\kbersatzschrift #1}\else{\kbersatzschrift #1}\fi\else ?\fi}
% Kopf- und Fußzeile (Befund 1, 2): Kopf Thema · Blattart, Fuß Kennung und Seite
\newcommand{\kbkopfzeile}[2]{\pagestyle{fancy}\fancyhf{}\renewcommand{\headrulewidth}{0pt}%
  \fancyhead[L]{\footnotesize\color{mbgrau}#1}%
  \fancyfoot[L]{\footnotesize\color{mbgrau}#2}%
  \fancyfoot[R]{\footnotesize\color{mbgrau}Seite \thepage}}
% Flattersatz überall (Befund 25)
\raggedright
% Titel des Blatts: Ich-kann-Satz, darunter Zweigzeile
\newcommand{\kbtitel}[2]{\par\noindent{\Large\bfseries\raggedright #1\par}%
  \ifblank{#2}{}{\par\smallskip\noindent{\small #2}\par}\medskip}
% Abschnitt: „Das kennst du schon“, „Zum Merken“
\newcommand{\kbabschnitt}[1]{\par\addvspace{12pt}\Needspace*{5\baselineskip}%
  \noindent{\bfseries #1}\par\nobreak\vspace{1pt}\noindent{\color{mbgrau}\rule{\linewidth}{0.4pt}}\par\nobreak}
% Hauptnummer: Titel hängt am ersten Block; jeder Block (Teilaufgabe) bricht
% nicht, passt er nicht mehr, rückt er ganz auf die nächste Seite (Befund 21).
\newif\ifkbtitel
\newsavebox{\kbbox}
\newlength{\kbhoehe}
\newlength{\kbnummerabstand}\setlength{\kbnummerabstand}{10pt}
\newlength{\kbteilabstand}\setlength{\kbteilabstand}{6pt}
\newenvironment{kbaufgabe}[1]{\stepcounter{aufgabe}\setcounter{teil}{0}%
  \gdef\kbtiteltext{\noindent\hangindent1.6em\hangafter1\makebox[1.6em][l]{\textbf{\theaufgabe.}}#1\par}%
  \global\kbtiteltrue\par\addvspace{\kbnummerabstand}}{\par}
\newenvironment{kbblock}{\par\addvspace{\kbteilabstand}\begin{lrbox}{\kbbox}%
  \begin{minipage}[t]{\linewidth}\raggedright\parskip\z@
  \ifkbtitel\kbtiteltext\vspace{2pt}\global\kbtitelfalse\fi
  \list{}{\leftmargin1.6em\labelwidth1.4em\labelsep0.2em\topsep\z@\partopsep\z@
    \parsep\z@\itemsep\z@\rightmargin\z@}\raggedright}%
  {\endlist\end{minipage}\end{lrbox}%
  \setlength{\kbhoehe}{\dimexpr\ht\kbbox+\dp\kbbox\relax}%
  \Needspace*{\kbhoehe}\noindent\usebox{\kbbox}\par}
% Paarsatz (Platz nutzen, Befund 27): zwei Hauptnummern oder zwei
% Teilaufgaben mit Zeichenaufgabe nebeneinander, als ein Block
\newcommand{\kbnummer}[1]{\stepcounter{aufgabe}\setcounter{teil}{0}%
  \noindent\hangindent1.6em\hangafter1\makebox[1.6em][l]{\textbf{\theaufgabe.}}#1\par\vspace{2pt}}
\newcommand{\kbhalb}[1]{\list{}{\leftmargin1.6em\labelwidth1.4em\labelsep0.2em\topsep\z@
  \partopsep\z@\parsep\z@\itemsep\z@}\raggedright #1\endlist}
\newcommand{\kbpaar}[2]{\par\addvspace{\kbnummerabstand}\begin{lrbox}{\kbbox}%
  \begin{minipage}[t]{0.485\linewidth}\raggedright\parskip\z@ #1\end{minipage}\hfill
  \begin{minipage}[t]{0.485\linewidth}\raggedright\parskip\z@ #2\end{minipage}\end{lrbox}%
  \setlength{\kbhoehe}{\dimexpr\ht\kbbox+\dp\kbbox\relax}\Needspace*{\kbhoehe}\noindent\usebox{\kbbox}\par}
\newcommand{\kbteilpaar}[2]{\item[]\noindent
  \begin{minipage}[t]{0.485\linewidth}\kbhalb{#1}\end{minipage}\hfill
  \begin{minipage}[t]{0.485\linewidth}\kbhalb{#2}\end{minipage}\par}
% Anweisung der Hauptnummer, eigene Zeile über der Teilaufgabe (Befund 5)
\newcommand{\kbanweisung}[1]{\item[]#1\par\vspace{2pt}}
% Kopfzeile einer Teilaufgabe: Buchstabe, Auftakt (halbfett oder Kontext),
% Prüfkennung klein rechts (Befund 3, 12, Beschluss c)
\newlength{\kbkennbreite}\setlength{\kbkennbreite}{1.9cm}
\newcommand{\kbkennung}[1]{{\footnotesize #1}}
\newcommand{\kbteil}[3]{\item[#1]\ifblank{#3}%
  {\parbox[t]{\linewidth}{\raggedright #2}}%
  {\parbox[t]{\dimexpr\linewidth-\kbkennbreite\relax}{\raggedright #2}\hfill
   \parbox[t]{\kbkennbreite}{\raggedleft\kbkennung{#3}}}\par}
% Frage in eigener Zeile (Befund 4), ein Satz je Zeile (Befund 10)
\newcommand{\kbfrage}[1]{\par\vspace{2pt}#1\par}
% Antwortfeld in eigener Zeile darunter, linksbündig (Befund 4, 8)
\newcommand{\kbantwort}[1]{\par\vspace{4pt}\noindent#1\par}
% Zwei Spalten: links Grafik/Tabelle/Aufgabe, rechts Antwort oder Rechenraum
% (Befund 27); oben bündig. Beginnt einen eigenen Listenpunkt ohne Marke.
\newcommand{\kbzweispaltig}[3][0.48]{\item[]\noindent
  \begin{minipage}[t]{#1\linewidth}\raggedright #2\end{minipage}\hfill
  \begin{minipage}[t]{\dimexpr0.98\linewidth-#1\linewidth\relax}\raggedright #3\end{minipage}\par}
% Linke Spalte mit Teilaufgabe (Buchstabe hängt links heraus wie in der Liste)
\newcommand{\kbspalte}[1]{\list{}{\leftmargin\z@\labelwidth1.4em\labelsep0.2em\topsep\z@
  \partopsep\z@\parsep\z@\itemsep\z@}\raggedright #1\endlist}
% Antwortfelder: 2,2 cm, in Punkten 1,2 cm; ohne Umbruch davor
\newcommand{\kbfeld}[1][]{\mbox{\underline{\hspace{2.2cm}}\ifblank{#1}{}{\,#1}}}
\newcommand{\kbfeldk}[1][]{\mbox{\underline{\hspace{1.2cm}}\ifblank{#1}{}{\,#1}}}
% Grafik oder Tabelle unter dem Text, linksbündig (Befund 8, 28)
\newcommand{\kbdarunter}[1]{\par\vspace{4pt}\noindent#1\par}
% Ankreuzen: je Option eine Zeile, linksbündig (Befund 6, 7)
\newcommand{\kbkreuz}[1]{\par\vspace{2pt}\noindent$\square$\ #1\par}
% Bearbeitungsraum (Befund 9): Schreibzeilen wie mathblatt (9 mm)
\newcommand{\kbraum}[1]{\schreibzeilen{#1}}
% Wertetabelle mit Köpfen im Textmodus (Befund 26): wie \wertetabelle der
% Vorlage, aber die Köpfe setzt der Aufruf selbst ($x$, „Zeit in h“).
\NewDocumentCommand{\kbwertetabelle}{o m m m}{%
  \def\mbkopf{}\def\mbwerte{}\def\mbleer{}\setcounter{mbn}{0}%
  \foreach \v in {#4}{\xappto\mbkopf{& $\v$}\xappto\mbleer{& }\stepcounter{mbn}\xdef\mbnn{\arabic{mbn}}}%
  \IfNoValueTF{#1}{}{\foreach \v in {#1}{\ifdefempty{\v}{\xappto\mbwerte{& }}{\xappto\mbwerte{& $\v$}}}}%
  \begingroup\renewcommand{\arraystretch}{1.3}%
  \edef\mbpre{|>{\noexpand\columncolor{mbkasten}}l|*{\mbnn}{>{\noexpand\centering\noexpand\arraybackslash}p{\noexpand\mbzell}|}}%
  \expandafter\mbtabstart\expandafter{\mbpre}\hline
  #2 \mbkopf \\ \hline
  \IfNoValueTF{#1}{\rule{0pt}{\mbzeile}#3 \mbleer \\ \hline}{#3 \mbwerte \\ \hline}
  \end{tabular}\endgroup}
% Merkkasten am Ende (Aufbau des Kompetenzblatts)
\newcommand{\kbmerk}[2]{\par\addvspace{14pt}\begin{lrbox}{\kbbox}\begin{minipage}[t]{\linewidth}%
  \noindent{\bfseries #1}\par\vspace{1pt}\noindent{\color{mbgrau}\rule{\linewidth}{0.4pt}}\par\vspace{4pt}%
  \uebersichtskasten{\raggedright #2}\end{minipage}\end{lrbox}%
  \setlength{\kbhoehe}{\dimexpr\ht\kbbox+\dp\kbbox\relax}\Needspace*{\kbhoehe}\noindent\usebox{\kbbox}\par}
% Lösungsblatt: Nummer und Buchstabe halbfett, Lösung daneben
\newcommand{\kbloesung}[2]{\par\addvspace{3pt}\noindent\hangindent2.8em\hangafter1\makebox[2.8em][l]{\textbf{#1}}#2\par}
\makeatother

\begin{document}
\kbkopfzeile{Potenz- und Exponentialfunktionen, Wachstum und Zerfall · Kompetenzblatt}{POT-K1}
% Kompetenzblatt POT-K1: potenz-exponentialfunktionen e2 „Wachstumstabelle fortschreiben“ (zusammenbau v0.7, Niveau FOR)
\kbtitel{Ich kann eine Wachstumstabelle mit dem Faktor fortschreiben.}{ab Kl. 10 · P10 ×5}
\setcounter{aufgabe}{0}

\kbabschnitt{Das kennst du schon}

% zone: potenz-exponentialfunktionen-zone-f7-v1
\begin{kbaufgabe}{Ich kann den Quotienten zweier Tabellenwerte berechnen.}
\begin{kbblock}
\kbteil{}{In einer Tabelle steht $45$, darunter $63$.}{}
\kbfrage{Berechne den Quotienten $63 : 45$.}
\kbantwort{\kbfeld}
\end{kbblock}
\end{kbaufgabe}

% zone: potenz-exponentialfunktionen-zone-f5-v1
\begin{kbaufgabe}{Ich kann den Prozentwert berechnen.}
\begin{kbblock}
\kbteil{}{Wie viel sind $20\,\%$ von $150\,€$?}{}
\kbantwort{\kbfeld[€]}
\end{kbblock}
\end{kbaufgabe}

% zone: potenz-exponentialfunktionen-zone-f1-v1
\begin{kbaufgabe}{Ich kann eine Erhöhung um p Prozent als Faktor rechnen.}
\begin{kbblock}
\kbteil{}{Ein Ticket kostet $50\,€$. Der Preis steigt um $10\,\%$.}{}
\kbfrage{Neuer Preis?}
\kbantwort{\kbfeld[€]}
\end{kbblock}
\end{kbaufgabe}

\kbabschnitt{Schritt für Schritt}

% leiter: potenz-exponentialfunktionen-e2-k2-s0-v1
\begin{kbaufgabe}{Ich kann erkennen, ob in einer Tabelle die Differenz oder der Quotient gleich bleibt.}
\begin{kbblock}
\kbanweisung{Kreuze an.}
\kbteil{}{Was bleibt in der Tabelle von Wert zu Wert gleich?}{}
\kbzweispaltig[0.36]{\kbwertetabelle[100,110,121]{$t$}{Wert}{0,1,2}}{\kbkreuz{die Differenz}\kbkreuz{der Quotient}}
\end{kbblock}
\end{kbaufgabe}

% leiter: potenz-exponentialfunktionen-e2-k2-s1-v1
\begin{kbaufgabe}{Ich kann den nächsten Wert mit dem Wachstumsfaktor berechnen.}
\begin{kbblock}
\kbteil{}{Die Anzahl der Mitglieder ist jetzt $250$ und wird jeden Schritt mit $1{,}15$ multipliziert.}{}
\kbfrage{Berechne den nächsten Wert.}
\kbantwort{\kbfeld}
\end{kbblock}
\end{kbaufgabe}

% leiter: potenz-exponentialfunktionen-e2-k2-s2-v1
\begin{kbaufgabe}{Ich kann den Startwert in die Tabelle eintragen.}
\begin{kbblock}
\kbteil{}{Eine Miete beträgt zu Beginn $300\,€$ und steigt jedes Jahr um $10\,\%$.}{}
\kbfrage{Ergänze das fehlende Feld.}
\kbdarunter{\kbwertetabelle[,330,363]{Jahr}{Miete in €}{0,1,2}}
\end{kbblock}
\end{kbaufgabe}

% leiter: potenz-exponentialfunktionen-e2-k2-s3-v1
\begin{kbaufgabe}{Ich kann eine Lücke in der Tabelle füllen.}
\begin{kbblock}
\kbteil{}{Der Bestand wächst jedes Jahr mit dem Faktor $1{,}1$.}{}
\kbfrage{Fülle die Lücke.}
\kbzweispaltig[0.45]{\kbwertetabelle[400,440,,532{,}4]{Jahr}{Anzahl}{0,1,2,3}}{\kbraum{2}}
\end{kbblock}
\end{kbaufgabe}

% leiter: potenz-exponentialfunktionen-e2-k2-s4-v1
\begin{kbaufgabe}{Ich kann mehrere Schritte auf einmal mit einer Potenz berechnen.}
\begin{kbblock}
\kbteil{}{Ein Bestand von $300$ wächst jedes Jahr mit dem Faktor $1{,}1$.}{}
\kbfrage{Berechne den Wert nach $4$ Jahren mit einer Potenz.}
\kbzweispaltig[0.46]{\kbantwort{\kbfeld}}{\kbraum{2}}
\end{kbblock}
\end{kbaufgabe}

% leiter: potenz-exponentialfunktionen-e2-k2-s5-v1
\begin{kbaufgabe}{Ich kann einen Schritt zurückrechnen.}
\begin{kbblock}
\kbteil{}{Im Jahr $3$ ist der Wert $363$, der Faktor ist $1{,}1$.}{}
\kbfrage{Berechne den Wert im Jahr $2$.}
\kbantwort{\kbfeld}
\end{kbblock}
\end{kbaufgabe}

% leiter: potenz-exponentialfunktionen-e2-k2-s6-v1
\begin{kbaufgabe}{Ich kann eine fehlende Zeitangabe in der Tabelle bestimmen.}
\begin{kbblock}
\kbteil{}{Der Wert nimmt jeden Tag mit dem Faktor $0{,}8$ ab.}{}
\kbfrage{Welche Zeitangabe fehlt im Kopf? \\ Prüfe mit dem Faktor.}
\kbzweispaltig[0.44]{\kbwertetabelle[500,400,256,163{,}84]{Tag}{Menge}{0,1,?,5}}{\kbraum{2}}
\end{kbblock}
\end{kbaufgabe}

% leiter: potenz-exponentialfunktionen-e2-k2-s7-v1
\begin{kbaufgabe}{Ich kann eine Abnahme mit einem Faktor kleiner als 1 fortschreiben.}
\begin{kbblock}
\kbteil{}{{\bfseries\boldmath Ein Stoff zerfällt}}{}
\kbfrage{Jede Stunde bleiben $70\,\%$. \\ Schreibe die Tabelle fort.}
\kbdarunter{\kbwertetabelle[1\,000,,,]{Stunde}{Menge in mg}{0,1,2,3}}
\end{kbblock}
\end{kbaufgabe}

% pruefung: potenz-exponentialfunktionen-e2-k2-s8-v1, potenz-exponentialfunktionen-e2-k2-s8-v3, potenz-exponentialfunktionen-e2-k2-s8-v7, potenz-exponentialfunktionen-e2-k2-s8-v9, potenz-exponentialfunktionen-e2-k2-s8-v5
\begin{kbaufgabe}{Ich kann das auch in Aufgaben aus der Prüfung.}
\begin{kbblock}
\kbteil{a)}{Ein Stromtarif kostet im ersten Jahr ($2026$) $80\,€$ im Monat und steigt jedes Jahr um $2{,}5\,\%$.}{P10 ’26 F}
\kbfrage{Ergänze die beiden fehlenden Felder.}
\kbzweispaltig[0.43]{\kbwertetabelle[,82{,}00,84{,}05,]{Jahr}{Euro}{2026,2027,2028,2029}}{\kbraum{2}}
\end{kbblock}
\begin{kbblock}
\kbteil{b)}{Der Stromverbrauch einer Firma lag $2015$ bei $36$ Tausend Kilowattstunden und steigt jährlich um $4\,\%$.}{P10 ’20}
\kbfrage{Ergänze die Tabelle.}
\kbzweispaltig[0.51]{\kbwertetabelle[,37{,}4,38{,}9,]{Jahr}{Tausend kWh}{2015,2016,2017,2018}}{\kbraum{2}}
\end{kbblock}
\begin{kbblock}
\kbteil{c)}{Ein Lamm wiegt bei der Geburt $4\,\mathrm{kg}$ und nimmt wöchentlich um $12\,\%$ zu.}{P10 ’19}
\kbfrage{Ergänze die drei fehlenden Werte (drei Nachkommastellen).}
\kbzweispaltig[0.51]{\kbwertetabelle[,4{,}480,,]{Woche}{Masse in kg}{0,1,2,5}}{\kbraum{2}}
\end{kbblock}
\begin{kbblock}
\kbteil{d)}{Der Wasserdruck in einer Leitung beträgt $800\,\mathrm{hPa}$ und sinkt je Kilometer um $12\,\%$.}{P10 ’17}
\kbfrage{Ergänze die beiden fehlenden Werte (ganze Zahlen).}
\kbdarunter{\kbwertetabelle[800,704,,545,]{Weg in km}{Druck in hPa}{0,1,2,3,6}}
\kbraum{2}
\end{kbblock}
\begin{kbblock}
\kbteil{e)}{Ein Kontrastmittel wird stündlich um $20\,\%$ abgebaut; zu Beginn sind es $8{,}00\,\mathrm{mg}$.}{P10 ’16}
\kbfrage{Ergänze den fehlenden Wert und die fehlende Zeitangabe.}
\kbdarunter{\kbwertetabelle[8{,}00,,5{,}12,3{,}28,2{,}62]{Zeit in h}{Menge in mg}{0,1,2,?,5}}
\kbraum{3}
\end{kbblock}
\end{kbaufgabe}

\kbmerk{Zum Merken}{Wachstumsfaktor: Zunahme um p Prozent heißt „mal q“ mit q = 1 + p/100, Abnahme um p Prozent heißt „mal q“ mit q = 1 − p/100: plus 6 \% → mal 1,06; minus 6 \% → mal 0,94. \\ Faktor aus der Tabelle: Wert geteilt durch den Wert davor; kommt immer derselbe Quotient heraus, ist der Faktor gefunden: 72 : 60 = 1,2. \\ Tabelle fortschreiben: nächster Wert = Wert davor mal q; ein Schritt zurück heißt geteilt durch q; der Anfangswert steht im Schritt null.}
\end{document}
